\section{Alternatives Considered}
\label{sec:alternatives}

Table~\ref{tab:alternatives} lists the alternatives encountered while building
the metamodels, with the reason each was rejected. The four discussed in prose
below are the ones where the rejection is not obvious.

\begin{table}[!htbp]
\centering
\small
\begin{tabular}{>{\raggedright\arraybackslash}p{0.30\linewidth}>{\raggedright\arraybackslash}p{0.62\linewidth}}
\hline
\textbf{Alternative} & \textbf{Why it was rejected} \\
\hline
Template the target files directly (Helm, Kustomize overlays) &
  No abstract model, so no constraint can be written and no fact can be derived
  once; the repetition the project exists to remove stays. \\
Annotate the target files with extra fields &
  Tried in the estate: the authoring shape and the resolved shape end up sharing
  one name, and no vocabulary can say which document is wrong. \\
Two stages, resolution private to the generator &
  The generator then takes decisions that appear in no model and can be
  reviewed nowhere. \\
A third metamodel for the generated files &
  The target syntax already has a schema; a third metamodel restates it. \\
ArchiMate as the source language &
  Covers the structural half and has no element for the operational half;
  expressing it means free-form properties, which gives up the closed
  vocabulary. \\
C4 model as the source language &
  Same gap, and it is a diagramming convention rather than a language with an
  abstract syntax to instantiate. \\
UML deployment or component diagram as the source language &
  An instance-level notation: it describes one arrangement, not the language
  every arrangement is written in. \\
A grammar first, with the metamodel inferred &
  Puts the concrete syntax in charge of the concepts; the same language has to
  be readable from two syntaxes. \\
Free-form strings instead of closed vocabularies &
  Tried in the estate: the one free-form field was wrong in every observed
  case. \\
An override mechanism for derived values &
  Hides a wrong derivation behind a per-process reason instead of exposing it
  for the whole class. \\
\hline
\end{tabular}
\caption{Alternatives Encountered and the Reason Each Was Rejected}
\label{tab:alternatives}
\end{table}

\subsection{ArchiMate and the C4 Model as the Source Language}

This is the alternative the Task 0 feedback proposed, and it was the most
seriously considered. ArchiMate is a full modelling language with a defined
metamodel, so unlike a diagramming convention it could in principle be
instantiated and transformed. The comparison in Section~\ref{sec:archimate}
shows why it is not used here: every structural concept maps cleanly, and every
operational concept maps to nothing.

The decisive point is what happens when the gap is filled. ArchiMate's
extension mechanism is properties: key-value pairs attached to elements. A
durability class expressed as a property is a string, so no constraint can
require that a process holding a recoverable volume also names what its data is,
and no transformation can enumerate the classes exhaustively. The project's
constraints are its main contribution, and they are written against closed
vocabularies. Adopting ArchiMate would mean writing those constraints against
property names instead, which is the free-form-string alternative in a different
form.

The C4 model is rejected for an additional reason: it is a set of diagram
conventions at four levels of abstraction, not a language with an abstract
syntax. There is no C4 metamodel to instantiate, so a C4 source language would
mean first inventing one.

Neither notation is discarded, though. Both remain good \emph{output}
notations, and Section~\ref{sec:archimate} records that an ArchiMate technology
view or a C4 deployment diagram is derivable from the resolved model by a second
model-to-text transformation, because the resolved model already contains every
fact such a view needs.

\subsection{A Third Metamodel for the Generated Files}

The Task 0 proposal carried three metamodels, one per stage of the pipeline.
It was reduced to two, and this is the clearest scope reduction the feedback
produced.

The argument for a third metamodel is that it would let the generation step be
checked structurally rather than as text. The argument against it is that the
target syntax is already schema-defined, so the third metamodel would be a
hand-maintained restatement of a schema that exists, with a drift problem of its
own and no constraint that the second metamodel cannot already express. The
generated files are therefore compared as text against reviewed expected output,
and the typed target resources live inside the second metamodel, where the
generation templates can walk them.

\subsection{Grammar First or Metamodel First}

The source language has a textual concrete syntax, so the language could have
been defined by a grammar with the metamodel inferred from it. It was defined
the other way round: the metamodel first, with the grammar as one syntax over
it.

The reason is that the same source model must be readable from two concrete
syntaxes, and the production implementation parses a different serialisation of
the same language. If the grammar were in charge, the two would agree only by
coincidence. With the metamodel in charge, both syntaxes are checked against one
structure, and that structure is exported as a machine-readable description that
both implementations are held to.

\subsection{One Resolved Model, or One Per Application}

The resolved model covers the whole estate in one document, rather than being
produced per application. Some of what it records cannot be computed for one
application alone: which shared edge carries a hostname, the order in which
units are applied, and what an application derives from the applications that
depend on \emph{it}. Producing a per-application model directly would either
recompute those properties inconsistently or omit them.

What an application's own repository receives is therefore a \emph{projection}:
the slice belonging to that application, obtained by filtering the estate-wide
model rather than by a separate computation, so the two cannot disagree about
what was decided.
